% palavras-chave
% iniciar todas com letras minúsculas, exceto no caso de abreviaturas

% Keywords precisam estar definidos ANTES do \maketitle

\begin{abstract}
Direct Digital Synthesis (DDS) is a signal synthesis method designed to generate arbitrary waveform from a single, fixed reference clock coupled with a memory device and a Digital-to-Analog Converter (DAC). This method can be used to synthesize basic sounds based on their recorded waveforms. Direct Digital Synthesis can be implemented not only in hardware but in software as well. The limitations being the memory and processing power of the processor involved in running the synthesis software.  Later with the use of electronic devices, these systems became smaller and more powerful, with many possibilities of configurations.
Microcontrollers are small computers that contain one or more Central Processing Units (CPU), memory and many different types of peripherals, all in the same package. There are many different types of microcontrollers architectures, for different applications and price ranges. Application-Specific Integrated Circuits (ASIC) designed to run a hardware DDS system are often expensive and not readily available. Microcontrollers on the other hand have become much more powerful, memory and performance wise, much more cheaper and available.
This work proposes a sound synthesis system running a software DDS architecture utilizing a pair of off-the-shelf, low cost microcontrollers, in order to make the design easy to replicate, to reprogram the firmware and to change the overall interface design and functionality.
\end{abstract}
